\documentclass[10pt,a4paper]{amsart}
\usepackage{amsmath,amssymb,amsthm,graphicx}
% For commutative diagrams we use Xy-pic
\usepackage[all]{xy}

% We use 2cell for 2-commutative diagrams.
\xyoption{2cell}
% Theorem environments.
%
\theoremstyle{plain}
\newtheorem{theorem}[subsection]{Theorem}
\newtheorem{proposition}[subsection]{Proposition}
\newtheorem{lemma}[subsection]{Lemma}

\theoremstyle{definition}
\newtheorem{definition}[subsection]{Definition}
\newtheorem{example}[subsection]{Example}
\newtheorem{exercise}[subsection]{Exercise}
\newtheorem{situation}[subsection]{Situation}

\theoremstyle{remark}
\newtheorem{remark}[subsection]{Remark}
\newtheorem{remarks}[subsection]{Remarks}

\numberwithin{equation}{subsection}

% Macros
%
\def\lim{\mathop{\mathrm{lim}}\nolimits}
\def\colim{\mathop{\mathrm{colim}}\nolimits}
\def\Spec{\mathop{\mathrm{Spec}}}
\def\Hom{\mathop{\mathrm{Hom}}\nolimits}
\def\Ext{\mathop{\mathrm{Ext}}\nolimits}
\def\SheafHom{\mathop{\mathcal{H}\!\mathit{om}}\nolimits}
\def\SheafExt{\mathop{\mathcal{E}\!\mathit{xt}}\nolimits}
\def\Sch{\mathit{Sch}}
\def\Mor{\mathop{\mathrm{Mor}}\nolimits}
\def\Ob{\mathop{\mathrm{Ob}}\nolimits}
\def\Sh{\mathop{\mathit{Sh}}\nolimits}
\def\NL{\mathop{N\!L}\nolimits}
\def\CH{\mathop{\mathrm{CH}}\nolimits}
\def\proetale{{pro\text{-}\acute{e}tale}}
\def\etale{{\acute{e}tale}}
\def\QCoh{\mathit{QCoh}}
\def\Ker{\mathop{\mathrm{Ker}}}
\def\Im{\mathop{\mathrm{Im}}}
\def\Coker{\mathop{\mathrm{Coker}}}
\def\Coim{\mathop{\mathrm{Coim}}}

% Boxtimes
%
\DeclareMathSymbol{\boxtimes}{\mathbin}{AMSa}{"02}

%
% Macros for moduli stacks/spaces
%
\def\QCohstack{\mathcal{QC}\!\mathit{oh}}
\def\Cohstack{\mathcal{C}\!\mathit{oh}}
\def\Spacesstack{\mathcal{S}\!\mathit{paces}}
\def\Quotfunctor{\mathrm{Quot}}
\def\Hilbfunctor{\mathrm{Hilb}}
\def\Curvesstack{\mathcal{C}\!\mathit{urves}}
\def\Polarizedstack{\mathcal{P}\!\mathit{olarized}}
\def\Complexesstack{\mathcal{C}\!\mathit{omplexes}}
% \Pic is the operator that assigns to X its picard group, usage \Pic(X)
% \Picardstack_{X/B} denotes the Picard stack of X over B
% \Picardfunctor_{X/B} denotes the Picard functor of X over B
\def\Pic{\mathop{\mathrm{Pic}}\nolimits}
\def\Picardstack{\mathcal{P}\!\mathit{ic}}
\def\Picardfunctor{\mathrm{Pic}}
\def\Deformationcategory{\mathcal{D}\!\mathit{ef}}
\usepackage{cleveref}
\crefname{originalproposition}{original proposition}{original propositions}
\Crefname{originalproposition}{Original proposition}{Original propositions}
\makeatletter
\newlabel{descent-proposition-same-cohomology-quasi-coherent}{{03DW}{source}{}{sourcepointer-1}{}}
\newlabel{descent-proposition-same-cohomology-quasi-coherent@cref}{{[originalproposition][0][]03DW}{[1][0][]source}{}{}{}}
\makeatother
\crefname{originalsection}{original section}{original sections}
\Crefname{originalsection}{Original section}{Original sections}
\makeatletter
\newlabel{flat-example-one-generator}{{0EVG}{source}{}{sourcepointer-2}{}}
\newlabel{flat-example-one-generator@cref}{{[originalsection][0][]0EVG}{[1][0][]source}{}{}{}}
\makeatother
\crefname{originalsection}{original section}{original sections}
\Crefname{originalsection}{Original section}{Original sections}
\makeatletter
\newlabel{flat-example-two-generators}{{0EVH}{source}{}{sourcepointer-3}{}}
\newlabel{flat-example-two-generators@cref}{{[originalsection][0][]0EVH}{[1][0][]source}{}{}{}}
\makeatother
\crefname{originallemma}{original lemma}{original lemmas}
\Crefname{originallemma}{Original lemma}{Original lemmas}
\makeatletter
\newlabel{flat-lemma-funny-blow-up}{{0EVJ}{source}{}{sourcepointer-4}{}}
\newlabel{flat-lemma-funny-blow-up@cref}{{[originallemma][0][]0EVJ}{[1][0][]source}{}{}{}}
\makeatother
\crefname{originallemma}{original lemma}{original lemmas}
\Crefname{originallemma}{Original lemma}{Original lemmas}
\makeatletter
\newlabel{flat-lemma-h-sheaf-colim-F}{{0EVK}{source}{}{sourcepointer-5}{}}
\newlabel{flat-lemma-h-sheaf-colim-F@cref}{{[originallemma][0][]0EVK}{[1][0][]source}{}{}{}}
\makeatother
\crefname{originalreference}{original reference}{original references}
\Crefname{originalreference}{Original reference}{Original references}
\makeatletter
\newlabel{lemma-refine-check-h}{{0EVI}{source}{}{sourcepointer-6}{}}
\newlabel{lemma-refine-check-h@cref}{{[originalreference][0][]0EVI}{[1][0][]source}{}{}{}}
\makeatother
\crefname{originallemma}{original lemma}{original lemmas}
\Crefname{originallemma}{Original lemma}{Original lemmas}
\makeatletter
\newlabel{sites-cohomology-lemma-colim-works-over-collection}{{0739}{source}{}{sourcepointer-7}{}}
\newlabel{sites-cohomology-lemma-colim-works-over-collection@cref}{{[originallemma][0][]0739}{[1][0][]source}{}{}{}}
\makeatother
\title{The perfection of the structure sheaf in characteristic p lies in the essential image of the h-topology pushforward}
\author{The Stacks Project Authors}
\date{stacks-project:etale-cohomology.tex:local-snapshot-20261011; GNU FDL 1.2 or later}
\begin{document}
\maketitle
\noindent\textit{Editorial note.} All delivered mathematical text is the author's own, copied from the pinned Stacks Project chapter file etale-cohomology.tex; no statement or proof step was rewritten, reordered or removed.\par\medskip
\noindent\textit{Editorial note.} The counted proof is the author's complete argument: the proof is written out in full for the configurations it treats, with the two cohomological inputs recorded at the start and the reduction to explicit short exact sequences of sections of O\textasciicircum{}\{perf\} in each case.\par\medskip
\noindent\textit{Editorial note.} The label lemma-h-sheaf-colim-F names the counted statement, and the proof environment immediately following it in the source is delivered verbatim as its proof.\par\medskip
\noindent\textit{Editorial note.} The counted statement is carried by the author's own lemma environment in the chapter source.\par\medskip
\noindent\textit{Editorial note.} The statement and proof are copied from the pinned chapter file; the seven cited inputs are kept as external pointers with their canonical Stacks tags, including the cross-chapter entries of the flatness chapter.\par\medskip
\noindent\textit{Editorial note.} The preamble is the chapter's own native preamble excerpt already used for this source: the Xy-pic setup (native preamble lines 14-18) and the native theorem declarations with the mathematical macros (native preamble lines 154-217).\par\medskip
\noindent\textit{Allowed standard prerequisites (editorial).}\begin{itemize}
\item The h-topology comparison morphism epsilon and its pushforward R epsilon\_*. Same chapter, the comparison between the h-topology and the fppf topology on (Sch/S), and the criterion for an object to lie in the essential image of R epsilon\_*. Supplies the target of the assertion and the criterion that the proof verifies.
\item Cohomology of the perfection as a filtered colimit. Cohomology on Sites, the lemma computing H\textasciicircum{}i of a filtered colimit of sheaves over a collection, applied to O\textasciicircum{}\{perf\} = colim\_F O. Supplies the identity H\textasciicircum{}i\_\{fppf\}(X, O\textasciicircum{}\{perf\}) = colim\_F H\textasciicircum{}i\_\{fppf\}(X, O) for quasi-compact quasi-separated X.
\item Almost blow-up squares and the exactness statements of the flatness chapter. More on Flatness, the one-generator and two-generator examples together with the funny blow-up and h-sheaf-colimit lemmas. Supplies the two configurations whose short exact sequences must be verified for the criterion.
\end{itemize}\medskip
\ifdefined\hypertarget\hypertarget{sourcepointer-1}{}\fi
\noindent\textit{Source reference.} Original proposition 03DW: sources/stacks-chow/tags.txt, source lines 4435--4435. Prerequisite pointer invoked at proof line(s) 23669; it names a standard fact of a sibling chapter that the counted argument uses, and every step of the counted proof built on it is retained verbatim in the delivered body.\par\medskip
\ifdefined\hypertarget\hypertarget{sourcepointer-2}{}\fi
\noindent\textit{Source reference.} Original section 0EVG: sources/stacks-chow/tags.txt, source lines 18586--18586. Prerequisite pointer invoked at proof line(s) 23673; it names a standard fact of a sibling chapter that the counted argument uses, and every step of the counted proof built on it is retained verbatim in the delivered body.\par\medskip
\ifdefined\hypertarget\hypertarget{sourcepointer-3}{}\fi
\noindent\textit{Source reference.} Original section 0EVH: sources/stacks-chow/tags.txt, source lines 18587--18587. Prerequisite pointer invoked at proof line(s) 23689; it names a standard fact of a sibling chapter that the counted argument uses, and every step of the counted proof built on it is retained verbatim in the delivered body.\par\medskip
\ifdefined\hypertarget\hypertarget{sourcepointer-4}{}\fi
\noindent\textit{Source reference.} Original lemma 0EVJ: sources/stacks-chow/tags.txt, source lines 18589--18589. Prerequisite pointer invoked at proof line(s) 23692; it names a standard fact of a sibling chapter that the counted argument uses, and every step of the counted proof built on it is retained verbatim in the delivered body.\par\medskip
\ifdefined\hypertarget\hypertarget{sourcepointer-5}{}\fi
\noindent\textit{Source reference.} Original lemma 0EVK: sources/stacks-chow/tags.txt, source lines 18590--18590. Prerequisite pointer invoked at proof line(s) 23685; it names a standard fact of a sibling chapter that the counted argument uses, and every step of the counted proof built on it is retained verbatim in the delivered body.\par\medskip
\ifdefined\hypertarget\hypertarget{sourcepointer-6}{}\fi
\noindent\textit{Source reference.} Original reference 0EVI: sources/stacks-etale-cohomology/pinned-original-source.tex, source lines 23614--23614. Prerequisite pointer invoked at proof line(s) 23659; it names a standard fact of the same chapter that the counted argument uses, and every step of the counted proof built on it is retained verbatim in the delivered body.\par\medskip
\ifdefined\hypertarget\hypertarget{sourcepointer-7}{}\fi
\noindent\textit{Source reference.} Original lemma 0739: sources/stacks-chow/tags.txt, source lines 9053--9053. Prerequisite pointer invoked at proof line(s) 23667; it names a standard fact of a sibling chapter that the counted argument uses, and every step of the counted proof built on it is retained verbatim in the delivered body.\par\medskip
\begin{lemma}
\label{lemma-h-sheaf-colim-F}
Let $p$ be a prime number. Let $S$ be a scheme over $\mathbf{F}_p$.
Consider the sheaf $\mathcal{O}^{perf} = \colim_F \mathcal{O}$
on $(\Sch/S)_{fppf}$. Then $\mathcal{O}^{perf}$ is in the essential
image of $R\epsilon_* : D((\Sch/S)_h) \to D((\Sch/S)_{fppf})$.
\end{lemma}

\begin{proof}
We prove this using the criterion of Lemma \ref{lemma-refine-check-h}.
Before check the conditions, we note that for a
quasi-compact and quasi-separated object $X$ of
$(\Sch/S)_{fppf}$ we have
$$
H^i_{fppf}(X, \mathcal{O}^{perf}) = \colim_F H^i_{fppf}(X, \mathcal{O})
$$
See Cohomology on Sites,
Lemma \ref{sites-cohomology-lemma-colim-works-over-collection}.
We will also use that $H^i_{fppf}(X, \mathcal{O}) = H^i(X, \mathcal{O})$, see
Descent, Proposition \ref{descent-proposition-same-cohomology-quasi-coherent}.

\medskip\noindent
Let $A, f, J$ be as in
More on Flatness, Example \ref{flat-example-one-generator}
and consider the associated almost blow up square.
Since $X$, $X'$, $Z$, $E$ are affine, we have no
higher cohomology of $\mathcal{O}$. Hence we only
have to check that
$$
0 \to
\mathcal{O}^{perf}(X) \to
\mathcal{O}^{perf}(X') \oplus \mathcal{O}^{perf}(Z) \to
\mathcal{O}^{perf}(E) \to 0
$$
is a short exact sequence. This was shown in (the proof of)
More on Flatness, Lemma \ref{flat-lemma-h-sheaf-colim-F}.

\medskip\noindent
Let $X, X', Z, E$ be as in
More on Flatness, Example \ref{flat-example-two-generators}.
Since $X$ and $Z$ are affine we have
$H^p(X, \mathcal{O}_X) = H^p(Z, \mathcal{O}_X) = 0$ for $p > 0$.
By More on Flatness, Lemma \ref{flat-lemma-funny-blow-up}
we have $H^p(X', \mathcal{O}_{X'}) = 0$ for $p > 0$.
Since $E = \mathbf{P}^1_Z$ and $Z$ is affine we also have
$H^p(E, \mathcal{O}_E) = 0$ for $p > 0$. As in the previous
paragraph we reduce to checking that
$$
0 \to
\mathcal{O}^{perf}(X) \to
\mathcal{O}^{perf}(X') \oplus \mathcal{O}^{perf}(Z) \to
\mathcal{O}^{perf}(E) \to 0
$$
is a short exact sequence. This was shown in (the proof of)
More on Flatness, Lemma \ref{flat-lemma-h-sheaf-colim-F}.
\end{proof}

\end{document}
